\subsection*{MT-M33. Finite Euler quotients disappear in a separated strip}

This is an unconditional obstruction to a bounded module
identification, not an obstruction to every possible
cutoff or geometric-kernel calculation.

\paragraph{The actual norm and the centered Euler operator.}
For a compact nonempty interval $I=[u,v]$ and $s\geq0$ put
\[
 H_{I,s}=L^2\left(\mathbf R,
       \max(e^{2ux},e^{2vx})(1+x^2)^{s/2}\,dx\right),
 \qquad (U_yf)(x)=f(x-y).
 \tag{M33.1}
\]
The compact norm-one group may be included with its
volume-one Haar measure; the same estimates hold.
Formula (M33.1) is exactly the fixed-character source norm
(9.2.2.31), original PDF p.~416.  The supremum over
subintervals in that definition is attained at the
endpoints of this compact $I$.

Let $T_af(x)=e^{ax}f(x)$.  The source completion
(10.1.1.3), original p.~434, is the closure of
$T_a\Theta$ in (M33.1), where $\Theta$ is its smooth
theta image.  Both original pages were rendered and
read.  The equivariance convention stated just above
that definition agrees with
\[
 U_yT_a=T_a(e^{-ay}U_y),\qquad
            T_aU_yT_a^{-1}=e^{ay}U_y.
 \tag{M33.2}
\]
These identities fix the sign of the centering twist.

In a compact character sector, write a uniformizer idele
as $(k_v,-h_v)$, with $h_v=\log q_v>0$, and put
$b_v=\beta(k_v)$, so $|b_v|=1$.
The uncentered difference of MT-M31 is
\[
 P_{v,0}=1-b_vU_{-h_v}.
\]
Transporting this literal operator by $T_a$ gives
\[
 P_{v,a}=T_aP_{v,0}T_a^{-1}
                 =1-b_ve^{-ah_v}U_{-h_v}.
 \tag{M33.3}
\]
Its Mellin multiplier is
$1-b_ve^{-(a+z)h_v}$, whose zeros have real part $-a$.
The operator is not the unscaled uncentered difference
when $a\ne0$.

\paragraph{The inverse with a quantitative convergent bound.}
If $a+u>0$, (M33.3) is a bounded automorphism of
$H_{I,s}$, with inverse
\[
 P_{v,a}^{-1}
       =\sum_{n=0}^{\infty}b_v^ne^{-anh_v}U_{-nh_v},
 \tag{M33.4}
\]
where convergence is in operator norm.  More precisely,
\[
 \|P_{v,a}^{-1}\|
 \leq 2^{s/4}\sum_{n=0}^{\infty}
       (1+n^2h_v^2)^{s/4}e^{-(a+u)nh_v}<\infty.
 \tag{M33.5}
\]
If $a+v<0$, the inverse instead is
\[
 P_{v,a}^{-1}
      =-\sum_{n=1}^{\infty}b_v^{-n}e^{anh_v}U_{nh_v},
 \tag{M33.6}
\]
with bound obtained by replacing the final exponential
in (M33.5) by $e^{(a+v)nh_v}$ and starting at $n=1$.
Thus every finite product
$P_{S,a}=\prod_{v\in S}P_{v,a}$ is a bounded
automorphism whenever $0\notin a+I$.

\begin{proof}
For every $y\in\mathbf R$, the endpoint exponential
weight and the polynomial inequality from MT-M31 give
\[
 \|U_y\|_{H_{I,s}\to H_{I,s}}
 \leq2^{s/4}(1+y^2)^{s/4}
                      e^{\max(uy,vy)}.
 \tag{M33.7}
\]
Indeed, after the change of variables $x=t+y$, each
endpoint weight acquires its factor $e^{2uy}$ or
$e^{2vy}$, and
$1+(t+y)^2\leq2(1+t^2)(1+y^2)$.
For $B_v=b_ve^{-ah_v}U_{-h_v}$ this proves
\[
 \|B_v^n\|\leq2^{s/4}(1+n^2h_v^2)^{s/4}
                                      e^{-(a+u)nh_v}.
\]
If $a+u>0$, the series is summable and $B_v^n\to0$
in operator norm.  The telescoping identity
$(1-B_v)\sum_{n=0}^{N}B_v^n=1-B_v^{N+1}$,
and its commuting counterpart, prove (M33.4).
If $a+v<0$, apply the same argument to
$B_v^{-1}=b_v^{-1}e^{ah_v}U_{h_v}$; the identity
$(1-B_v)^{-1}=-\sum_{n\geq1}B_v^{-n}$ proves
(M33.6).  The factors and their inverses commute,
so their finite products are automorphisms too.
\end{proof}

\paragraph{All closed invariant subquotients have the same property.}
Let $Y_0\subset Y\subset H_{I,s}$ be closed subspaces
invariant under every $U_y$, $y\in\mathbf R$.
Then $P_{S,a}$ and its inverse preserve $Y_0$ and $Y$:
each term of (M33.4) or (M33.6) preserves them and
they are closed under the norm limit.
Consequently $P_{S,a}$ is an automorphism on $Y/Y_0$
with its quotient norm.  It is enough to assume
invariance under the particular shifts appearing in
the inverse series and their products.

In the actual arithmetic setting,
first work in one fixed compact character sector and put
\[
          J_{a,I,s}=\overline{T_a\Theta}^{\,H_{I,s}}
 \tag{M33.8}
\]
This space is invariant under all norm translations by (M33.2),
since $\Theta$ is translation invariant.
Therefore, without GRH,
\[
 P_{S,a}J_{a,I,s}=J_{a,I,s},\qquad
 \overline{T_aP_{S,0}\Theta}^{\,H_{I,s}}=J_{a,I,s}
                   \quad(0\notin a+I).
 \tag{M33.9}
\]
For the second equality, use
$T_aP_{S,0}=P_{S,a}T_a$, density in (M33.8), and
the fact that an automorphism maps a dense subspace
onto a dense subspace of its whole target.
In particular
\[
       J_{a,I,s}/\overline{T_aP_{S,0}\Theta}=0.
 \tag{M33.10}
\]
For $a=1/2$ and $I=[-\alpha,\alpha]$ this applies
to every $0\leq\alpha<1/2$, including the line
$\alpha=0$.  This concerns the finite Euler quotient,
not the quotient by the full global theta ideal:
the latter need not vanish.

The full compact-group assertion has the same proof:
replace each scalar $b_v$ by the corresponding compact
unitary translation.  Its integer powers have norm one,
so the bounds (M33.5)--(M33.7) are uniform in the compact
character.  For that version the closed subspaces must
also be invariant under these compact translations and
their inverses.  The full theta closure has this
invariance.  Norm-translation invariance alone is the
hypothesis for the fixed-character statement above.

\paragraph{The exact intertwining obstruction.}
Let $Z_{S,s}$ be the corrected nonzero local quotient of
MT-M32, in the original unreflected coordinates:
\[
       Z_{S,s}=H_s/\overline{P_{S,0}H_s}.
 \tag{M33.11}
\]
For instance one prime and $s>1$ gives the nonzero
scalar circle model (M32.14c).
Ordinary norm translations act on this quotient.
Give the same underlying Hilbert space its centered
action
\[
                   V_y=e^{-ay}U_y.
 \tag{M33.12}
\]
Then the polynomial in these operators
\[
       \prod_{v\in S}(1-b_ve^{-ah_v}V_{-h_v})
       =\prod_{v\in S}(1-b_vU_{-h_v})
       =0\quad\hbox{on }Z_{S,s}.
 \tag{M33.13}
\]
The last equality follows from its definition as a
quotient by the closure of the range, and does not
require the actual range to be closed.

Assume $0\notin a+I$, and let $W=Y/Y_0$ be any
closed invariant subquotient as above, with its
inherited norm-translation action.  Every bounded
linear map
\[
       T:(Z_{S,s},V)\longrightarrow W
       \quad\hbox{or}\quad
       T:W\longrightarrow(Z_{S,s},V)
 \tag{M33.14}
\]
which intertwines each norm translation is zero.
For the first direction, applying the same operator
polynomial to the intertwining relation gives
$P_{S,a}T=0$, while $P_{S,a}$ is injective on $W$.
For the second direction it gives $TP_{S,a}=0$,
while $P_{S,a}$ is surjective on $W$.
This proves both assertions.  The twist in (M33.12)
is essential: comparing the untwisted circle action
would compare a different polynomial and would not
prove this statement.

\paragraph{Location of the remaining source comparison.}
The source itself separates the spaces before attempting
the later comparison.  Theorem 11.1.2.1, p.~480,
is explicitly a theorem on the $0$-line.  The local
orbital quotient (11.2.2.1), p.~489, is also explicitly
defined with $a=0$, and Theorem 11.2.2.1,
pp.~490--491, gives its local trace.
In contrast, the absolute and invariant Paley--Wiener
constructions at pp.~487--489 allow a general center
$a$, while Theorem 11.3.2.1, pp.~495--496, sets
$a=1/2$ and permits $I=[-\alpha,\alpha]$ for
$\alpha\geq0$.  Its proof invokes the zero-line
carving theorem and the local orbital trace after
cutoffs and shell degeneration.  The exact norm and
centering convention at pp.~416 and 434 have been
visually checked for this comparison; the cited
transition statements were read in the supplied text.

Equations (M33.9)--(M33.14) show that the needed transition
cannot be a nonzero bounded identification of the
centered local circle module with a closed invariant
subquotient of the narrow Hilbert space.  Any valid
geometric calculation must instead justify the actual
non-equivariant cutoff operators and their limits.
Neither the local Poisson identity in MT-M32 nor the
conditional full narrow-strip spectral trace in MT-M28
supplies that justification.  No conclusion that the
Hecke-GRH endpoint is false is drawn here.

In the cote 1 trace programme (scans 79--80 and 111--112),
this locates the issue in transport from a local
closed-point Euler module to the global trace space.
The local Poisson calculation survives; the separated
strip changes the relevant translation spectrum, so
the geometric comparison needs its own cutoff-limit
argument.
